\section{Introduction}\label{sec:intro}

\HB{
This article presents an arbitrage-free market model for for smooth surfaces of European option prices.

The core architecture follows the SANOS approach~\cite{SANOS2026} which models the stock price as the product
of a time-and-state discrete process~$X$ and a smoothing process~$Y$, in our case Black~\&~Scholes.
The former is used to encode the option market using a discrete form of local volatilities as code book. That does not mean
that the model is a local volatility model -- we merely use the (discrete) local volatilities as a convenient representation of the option surface.
Compared to using implied volatilities the advantage is that mere positivity of the local volatilities is sufficient to ensure
the the option market is free of static arbitrage at all times. We have shown in~\cite{SANOS2026} that this representation
can be fitted to real life market data very efficiently using linear programming.

In~\cite{DYSANOS2026} we have shown how to estimate the dynamics such (discrete) local volatilities of such a model under the
statistical measure~$\mathbb{P}$ which fit very well to real-life observed option bid/asks.

The key
contribution in this paper is showing how to construct a corresponding risk-neutral measure~$\mathbb{Q}$ or, equivalently,
the risk-neutral drift of the model. Our approach is closely related to the seminal but unpublished work~\cite{Wissel2007}.
}

An option market model reverses the usual direction of construction: instead of
specifying a process for the underlying and deriving option prices from it, one
models the code-book of prices directly, so that a calibration is a change of
coordinates rather than an inverse problem. Two conditions must then be
maintained at once. The state must remain free of \emph{static} arbitrage, and
its dynamics must admit one equivalent local martingale measure common to the
underlying and to every traded option. The second is the option analogue of an
HJM drift restriction: the volatility structure may be specified freely, and the
drift is then determined. The lattice here is \emph{finite}, finitely many
strikes and finitely many expiry dates, which is what is traded.

\paragraph{What is already known.} Schweizer and Wissel treat the problem in the
continuum in two companion papers, one in each direction of the surface:
\cite{Schweizer2007} varies maturity at a single strike, and
\cite{SchweizerWissel2008} varies strike at a single maturity. In each the
code-book is chosen so that its positivity encodes absence of arbitrage, the
drift restriction is derived, and existence is proved. The multi-strike
multi-maturity case, which is ours, is stated as open in
\cite{SchweizerWissel2008}, with the expectation that a new parametrisation will
be needed. Carmona and Nadtochiy
\cite{CarmonaNadtochiy2008} give a continuum local-volatility code-book with a
drift restriction, but no existence theorem. On a finite lattice, and answering
the multi-maturity question, stands Wissel \cite{Wissel2007}; his construction is
the one this paper is closest to, and it is described separately below.

\paragraph{Modern approaches.} The question was reopened from the data side, by
the scarcity of option-market histories for training and evaluating hedging
strategies \cite{BuehlerEtAl2019DeepHedging}. Cohen, Reisinger and Wang
\cite{CohenReisingerWang2021} build neural stochastic differential equations on a
floating lattice and penalise a least-squares residual in place of the drift
restriction.
\cite{WieseEtAl2019OptionMarkets} simulate an equity option book with a
generative adversarial network whose state is the discrete local variance grid
compressed by principal components, so that static arbitrage is absent by
positivity of that grid; that is the code-book choice this paper inherits, and
\cite{WieseEtAl2021MultiAsset} extend it to a multi-asset universe.
\cite{DYSANOS2026} decode statically arbitrage-free surfaces from a
low-dimensional code, fit a linear state process to it, and test for dynamic
arbitrage after the fact, reporting it in $12.5\%$ of states at zero transaction
cost.

None of the four enforces the drift restriction exactly, and that is where this paper starts:
the cross-sectional constraints govern each date's surface and its terminal
payoff cone, and leave the transition of a fixed contract between two dates
unconstrained. The choice of grid is part of the same point. All four index
their nodes relative to the spot rather than in fixed cash terms, and on such a
grid a node is not one contract through time, so there is no single instrument
whose price is required to be a martingale and nothing for an HJM-type
restriction to constrain. Fixing the strikes in cash is what buys the
restriction here.

\paragraph{The closest prior work.} Wissel \cite{Wissel2007} works on a finite lattice of fixed cash strikes and
fixed expiry dates, parametrises it by a price level together with positive local
implied volatilities, derives the no-arbitrage drift restrictions, and proves
constructive strong existence and uniqueness for a class of coefficients. Behind
the front expiry his local implied variance is \emph{exactly} the backward
discrete local variance used here, an algebraic identity on the grid, and his
drift restriction is, after a change
of measure, exactly the one we state in Theorem~\ref{thm:drift}. We claim no
priority for it; see Appendix~\ref{app:wissel}, which gives the identification in
full. That working paper appears never to have been published in archival
form.

\subsection*{What this paper does}

We keep the same lattice and change the chart: the state is the vector of
gridded \emph{densities} together with the discrete local variances, rather than
prices together with quantiles of digitals. Three consequences follow, and they
are the contributions.

\begin{enumerate}[leftmargin=1.8em,itemsep=0.45em]
\item \textbf{The spot is a read-off, and it settles exactly.}
The spot is $S_t=K^{\!\top}q^j_t$ at every expiry $j$ (Proposition~\ref{prop:spot}),
an identity of the chart holding at all $t$. At the front
expiry Proposition~\ref{prop:collapse} exhibits a family of front volatilities
under which $q^1$ collapses onto a single node, so the spot lands on a strike
and every gridded contract settles at its exact payoff
(Section~\ref{sec:nonpinning}). That pinning is a property of the volatility
chosen and not of the chart: by Remark~\ref{rem:bary} the terminal states
the chart admits carry every spot in the grid range, whereas in
\cite{Wissel2007} the pinning follows from the chart whatever the coefficients
(Remark~\ref{rem:pinning}). We do not construct dynamics reaching an off-lattice
terminal spot; that is left as future work.

\item \textbf{Existence, with no condition on the wings.} For every bounded and
locally Lipschitz volatility structure $b$ and any number of driving factors, the
model has a unique strong solution on $[0,T]$ for every $T<T_1$, positive by
construction; it extends to $T_1$ (Remark~\ref{rem:horizon}) and, restarting at each
expiry with the settling front, to the whole ladder $[0,T_M]$, unique given that
front (Corollary~\ref{cor:patch}). The quantity that must be controlled is
$\eta=\psi/q$, the density volatility relative to the density, and its growth across
expiries is bounded by the grid alone: Theorem~\ref{thm:wing} gives $\Theta_j\le N+1$
uniformly in the state and in the volatility structure. \emph{Wissel \cite{Wissel2007} must assume a bound of this
kind; here it is a theorem.} The front expiry costs one hypothesis, and it is a mild one: $\eta^1$
must be integrable in $t$, not bounded, because $\eta^j$ is affine in $\eta^1$
and an It\^o equation needs an integrable drift rather than a bounded one. The
front volatility used here diverges at rate $(T_1-t)^{-p}$, and the two demands
made of $p$ leave a window: $p<1$ is what Assumption~\ref{ass:F} asks, while
$p\ge\tfrac12$ is what makes the front density collapse, so that every gridded
contract settles at its exact payoff (Proposition~\ref{prop:collapse}).

\item \textbf{The drift is pointwise, and the count works out.} In density
coordinates the drift equation is $N+2$ conditions on $N$ unknowns, and is
consistent because the left null space of $\widehat G_j$ is exactly
$\mathrm{span}\{\one,K\}$, conservation of mass and of the spot
(Lemma~\ref{lem:null}). The state-independent factor then cancels from both
sides and the drift is available nodewise,
$\alpha^{j,i}=-\ip{\eta^{j,i}}{\beta^{j,i}}$, with no matrix inversion and no
rank condition.
\end{enumerate}

\subsection*{Scope, and what is not claimed}

The over-determination of the third contribution is not intrinsic: it is created
by the choice of density coordinates and removed by the same choice, and in
price coordinates the system is square and the question does not arise.

The existence statement is for a fixed finite grid. The constant is
$(M-1)(N+1)$ and depends on the grid. The one
hypothesis, Assumption~\ref{ass:F}, is a condition on the front volatility and
not on the state, and it is sufficient rather than necessary.
